% from main.tex
%%%%%%%%%%  scan idx 276  |  folio 263  %%%%%%%%%%
résultent de (5.1) appliqué au topos «~ponctuel~» annelé par \(\Gamma(X,\mathcal O_X)\) ; il en est de même, en vertu de (II 2.2.10 a)), de l'implication \(\mathrm b_{\qcoh})\Longrightarrow\mathrm b'_{\qcoh})\) pour \(E\) pseudo-cohérent. Enfin l'équivalence \(\mathrm a)\Longleftrightarrow\mathrm a_{\qcoh})\), ou, ce qui revient au même, \(\mathrm a')\Longleftrightarrow\mathrm a'_{\qcoh})\) découle de la dernière remarque de (ii) (et du fait que les faisceaux quasi-cohérents «~se prolongent~», \(X\) étant quasi-séparé).

\textbf{Proposition 5.3.} (propriété d'échange) (EGA III 7.7.5 II). Soient \(f:X\to Y\) un morphisme de schémas quasi-compacts et quasi-séparés, \(E\in\Ob D^b(X)\) un complexe à cohomologie quasi-cohérente et d'amplitude plate finie rel. à \(f\), et \(p\in\mathbb Z\). Considérons les conditions suivantes :
\begin{enumerate}[label=\alph*)]
\item[a)] Le foncteur \(M\mapsto R^pf_*(E\otimes_Y^{\mathbf L}M)\) de \(\Qcoh(Y)\)\footnote{\(\Qcoh(Y)\) = catégorie des \(\mathcal O_Y\)-Modules quasi-cohérents.} dans \(\Qcoh(Y)\)\footnote{en vertu de (3.7).} est exact à droite.
\item[a')] La flèche canonique
\[
 R^pf_*(E)\otimes M\longrightarrow R^pf_*(E\otimes_Y^{\mathbf L}M)
\]
(composée de la flèche canonique \(R^pf_*(E)\otimes M\longrightarrow H^p(Rf_*(E)\otimes^{\mathbf L}M)\) et de l'isomorphisme (3.7) \(H^p(Rf_*(E)\otimes^{\mathbf L}M)\xrightarrow{\sim}R^pf_*(E\otimes_Y^{\mathbf L}M)\)) est un isomorphisme pour tout \(M\in\Ob\Qcoh(Y)\).
\item[b)] Le foncteur \(M\mapsto R^{p+1}f_*(E\otimes_Y^{\mathbf L}M)\) de \(\Qcoh(Y)\) dans \(\Qcoh(Y)\) est exact à gauche.
\item[b')] La flèche canonique
\[
 R^{p+1}f_*(E\otimes_Y^{\mathbf L}M)\longrightarrow\underline{\mathrm{Hom}}(H^{-p-1}(Rf_*(E)^{\vee}),M)
\]
(composée de l'isomorphisme (3.7) \(R^{p+1}f_*(E\otimes_Y^{\mathbf L}M)\xrightarrow{\sim}H^{p+1}(Rf_*(E)\otimes^{\mathbf L}M)\) et de la flèche canonique (5.1 (\(*\))) \(H^{p+1}(Rf_*(E)\otimes^{\mathbf L}M)\longrightarrow\underline{\mathrm{Hom}}(H^{-p-1}(Rf_*(E)^{\vee}),M)\)) est un isomorphisme pour tout \(M\in\Ob\Qcoh(Y)\).
\item[c)]
